% ======================================================================
% 06-results.tex — results that survive re-derivation
% ======================================================================
\section{Results that survive re-derivation}
\label{sec:results}

The point of the mechanisms is not tidiness. It is that some claims now survive
attempts to break them, and the ones that do not are labelled. This section
reports the survivors and the caveats together; the full tables are in
Appendix~\ref{app:tables}.

\subsection{One dimension: the baseline is honest, the FNO is fixed}
\label{sec:res1d}

Figure~\ref{fig:1d} shows the PINN against the Ogata--Banks reference. The
qualitative agreement is not the result; the result is the table.

\begin{figure}[htbp]
\centering
\includegraphics[width=0.48\textwidth]{figs/field.png}\hfill
\includegraphics[width=0.48\textwidth]{figs/profiles.png}
\caption{1D sandbox: concentration field and outlet profiles for the hard-BC
PINN against the analytical reference. Generated by \code{SCRIPTS/hydrology.py}
and \code{SCRIPTS/solve\_1d\_pinn.py}.}
\label{fig:1d}
\end{figure}

\begin{table}[htbp]
\centering
\small
\begin{tabular}{lccccc}
\toprule
 & median & IQR & min & max & spread \\
\midrule
$E_{\mathrm{rel}}$ (full grid) & $1.61\times10^{-2}$ & $2.18\times10^{-2}$ & $1.17\times10^{-2}$ & $8.14\times10^{-2}$ & $7.0\times$ \\
$E_{\mathrm{rel}}^{\mathrm{int}}$ ($\xi,\tau>0.05$) & $1.42\times10^{-2}$ & $2.59\times10^{-2}$ & $6.64\times10^{-3}$ & $8.16\times10^{-2}$ & $12.3\times$ \\
\bottomrule
\end{tabular}
\caption{Phase 1 hard-BC baseline over five seeds. The distribution is
right-skewed (mean $3.19\times10^{-2}$, twice the median), so median and IQR are
reported. Artifact: \code{runs/sweeps/baseline\_n5}; harness:
\code{SCRIPTS/seed\_sweep.py}.}
\label{tab:1d}
\end{table}

Two corrections follow from Table~\ref{tab:1d}. First, the corner-excluded
metric does not rescue the target: the best seed
($6.64\times10^{-3}$) is below $10^{-2}$, the median ($1.42\times10^{-2}$) is
not. Second, the corner is not the governing artefact: the two metrics are
rank-correlated (Spearman $\rho=1$), and their ratio ranges from $0.57$ to
$1.00$, so excluding the corner never reorders the runs.

The FNO results are reported with the same discipline. The original operator
scored $0.973$ (worse than zeros); with coordinate channels and 250 epochs it
scored $0.0146$ on one seed. At 120 epochs across five seeds the median is
$3.02\times10^{-2}$ (IQR $3.2\times10^{-3}$, range $2.35\times10^{-2}$ to
$5.98\times10^{-2}$). The honest 120-epoch figure is therefore about 3\%, not
1.5\%; the 250-epoch number still needs its own multi-seed median before it is
quoted as a result. Mode truncation is monotone (modes 4/12/24 $\to$
2.79/2.35/2.19\%, single seed), and resolution invariance degrades gracefully
(train 64$^2$ $\to$ 4.14\%, train 32$^2$ $\to$ 6.92\%, evaluated at
128$^2$), which is the behaviour a genuine operator should show.

\subsection{Two dimensions: two policies, one certificate}
\label{sec:res2d}
\label{sec:res3dcert2d}

The $160\times96$ benchmark is a $200\times120$\,m aquifer, log-normal
conductivity ($\sigma=0.65$, $\ell=15$\,m, $n_{KL}=1529$), nine candidate wells,
$M_{\text{train}}=100$ and $M_{\text{test}}=40$. Two runs used the same search
and differ only in floating-point precision.

\begin{table}[htbp]
\centering
\small
\caption{Two independently certified optima of the same benchmark. Both satisfy
every gate; they differ by 1.2\% in cost while using different well subsets.}
\label{tab:2d}
\begin{tabular}{lrr}
\toprule
 & float64 & float32 (mixed) \\
\midrule
$\mathbf{q}^\star$ [$\mathrm{m^3/d}$] & $(0,9.72,5.86,0,3.38,0,0,0,0)$ & $(0,9.33,8.16,0,0,0,0,0,0)$ \\
total pumping & $19.0$ & $17.5$ \\
active wells & $3$ & $2$ \\
$J$ (train / test) & $433{,}945$ / $434{,}008$ & $428{,}592$ / $428{,}579$ \\
$\Cvar_{5\%}(G)$ (train / test) & $+0.018$ / $-0.007$ & $+0.005$ / $+0.009$ \\
reliability (train / test) & $0.990$ / $0.975$ & $0.990$ / $0.975$ \\
holdout during search & $1.000$ & $0.975$ \\
max.\ mass-balance error & $4.2\times10^{-11}$ & $2.3\times10^{-5}$ \\
wall clock (T4) & $4.14$\,h & $1.67$\,h \\
\bottomrule
\end{tabular}
\end{table}

Two facts matter more than the numbers. Both optima sit \emph{on} the
constraint ($\Cvar\approx0$ from either side), which is what the solution of a
chance-constrained program must do: they are the cheapest policies that are
still safe, not policies that happen to be safe. And the two policies differ by
1.2\% in cost while using different well subsets, so near the optimum the
objective is flat along ``which wells carry the duty''. Replicas with offset
scenario seeds reproduce the same policy digit for digit, which is why the flat
direction is a property of the cost model, not of sampling noise.

\begin{table}[htbp]
\centering
\small
\caption{Out-of-sample certificate on 400 scenarios with 95\% Wilson intervals.
Both clear the $1-\alpha-1\%=0.94$ gate, but not equally. A fresh draw from a
seed no run saw leaves the cost intact and moves the tail.}
\label{tab:cert}
\begin{tabular}{lrrr}
\toprule
policy & violations & reliability & 95\% Wilson CI \\
\midrule
float64 $\mathbf{q}^\star$ & $1/400$ & $0.998$ & $[0.986,\ 1.000]$ \\
float32 $\mathbf{q}^\star$ & $7/400$ & $0.983$ & $[0.964,\ 0.991]$ \\
\bottomrule
\end{tabular}
\end{table}

Table~\ref{tab:cert} is the strongest evidence in the article because it is the
one that does not come from the run's own summary. On scenarios drawn from seed
900000, which no run had seen, the float32 policy gives $J=428{,}509$ against
the run's $428{,}592$ (0.02\% apart), reliability $391/400=0.978$ with Wilson
$[0.958,0.988]$, and $\Cvar_{5\%}(G)=+0.616$ instead of $+0.005$. The policy is
safe; the tail it reported is a property of the sample it optimized on. The same
draw at $M=100$ gives $4/100=0.960$ with a Wilson interval of $[0.902,0.984]$,
which is inconclusive: the search's own sample size cannot certify the policy it
produced. Figure~\ref{fig:certificate} makes that argument visual: the intervals
only resolve the two policies at $M=400$.

\begin{figure}[htbp]
\centering
\includegraphics[width=0.58\textwidth]{figs/reliability_certificate.png}
\caption{Empirical reliability with its 95\% Wilson interval as a function of the
number of out-of-sample scenarios, for both certified policies (drawn on the
host run's test seeds; the seed-900000 draw above is the independent check). At
$M\le100$ the intervals overlap, straddle the $0.94$ gate and cannot separate the
two policies; only at $M=400$ does the weaker one ($0.983$) fall clearly below the
stronger ($0.998$). A point estimate of $0.975$ at $M=40$ is consistent with a
true reliability anywhere in $[0.87,1.00]$, which is why the certificate is
quoted as an interval. Generated by \code{SCRIPTS/make\_figures.py} from the
\code{check\_tail.py --ladder} output.}
\label{fig:certificate}
\end{figure}

\paragraph{A cheaper policy that fails the audit.} The certificate above is
built from policies a search \emph{accepted}. The complementary case is a
deliberately weaker search: the same objective at $160\times96$ with
$M_{\text{train}}=M_{\text{test}}=M_{\text{holdout}}=40$, two starts, 400
evaluations per start, float32 with the float64 Krylov recurrence, run twice on
different random seeds (7 and 8) so that only the starting points differ.
Table~\ref{tab:restart} is the result.

\begin{table}[htbp]
\centering
\small
\caption{Restarts do not find a second optimum at $M_{\text{train}}=40$; they
find the same fragile one twice. Both replicas optimise the same 40-scenario
objective, so the only difference between them is the starting points.}
\label{tab:restart}
\begin{tabular}{lrrrrl}
\toprule
replica & start 1 (det.) & start 2 (random) & $\mathbf{q}^\star$ & train / test rel. & verdict \\
\midrule
seed 7 & $409{,}863$ & $421{,}639$ & $[2.59,0,9.72,0\ldots]$ & $0.975$ / $0.900$ & fails \\
seed 8 & $409{,}863$ & $415{,}576$ & $[2.59,0,9.72,0\ldots]$ & $0.975$ / $0.900$ & fails \\
\bottomrule
\end{tabular}
\end{table}

Both random starts converged \emph{worse} than the deterministic one, so both
replicas returned the identical $\mathbf{q}^\star$ ($12.32\ \mathrm{m^3/d}$ from
two wells): at this sample size the landscape has one dominant attractor, and it
is the fragile policy. On the seed-900000 draw the same audit that passed the two
certified policies reports $J=409{,}645$---$4.4\%$ cheaper than the accepted
policy---with $34/400$ violations, reliability $0.915$ and Wilson
$[0.884,0.939]$, an interval entirely below the gate. The violation count grows
with the sample ($0/16$, $1/20$, $2/40$, $3/60$, $5/100$, $6/150$, $13/200$,
$28/300$, $34/400$), the opposite of the accepted policies, whose count stays
flat while the tail mean drifts; under-sampling can hide violations or exaggerate
them, so the sign is not predictable and the sample cannot be sized by watching
$\Cvar$.

The failure is statistical, not numerical. The 802 evaluations produced no
non-finite objective and no stalled solve, so float32 at $M=40$ did its job and
the sample size did not. This is also the first live catch for the mid-run
guard: a disjoint holdout archive, re-evaluated every eight sweeps at about $1\%$
of a sweep's cost, reported $\text{rel}=0.900$ and $\Cvar=+1.63$ at iteration 16
on \emph{both} replicas, and again at iteration 24---turning a post-mortem into an
abort decision while the VM was still running.

Figure~\ref{fig:plume} shows what the policy does to the plume. On 80 held-out
scenarios the float64 policy lowers the peak concentration from $35.81$ to
$0.433\ \mathrm{mg/L}$ (median, IQR $-97.4\%$ to $-99.3\%$) and the
compliance-point concentration from $17.23$ to $0.0260\ \mathrm{mg/L}$, raising
the number of scenarios meeting the $1.0\ \mathrm{mg/L}$ limit from $1/80$ to
$78/80$. The plume is displaced rather than erased: its peak moves
down-gradient to $x\approx115$\,m, adjacent to the active well.
Figure~\ref{fig:chance} shows the tail of the violation distribution that these
medians average over.

\begin{figure}[htbp]
\centering
\includegraphics[width=0.88\textwidth]{figs/contamination.png}
\caption{Untreated contamination and the certified policy on the 2D benchmark.
Each panel carries its own colour scale so the cleaned field is visible; active
wells are annotated with their rates. Artifact:
\code{out/certified\_20261120-111849/figures/contamination.png}.}
\label{fig:plume}
\end{figure}

\begin{figure}[htbp]
\centering
\includegraphics[width=0.92\textwidth]{figs/chance_tail.png}
\caption{The constraint behind the certificate, at $M=40$ out-of-sample
scenarios and the float64 policy. (a) The distribution of the worst
compliance-point violation $G$: 39 scenarios sit at the $\sim-1$ floor because the
plume is fully captured, one violates. The priced $\Cvar_{5\%}(G)\approx0$ -- the
policy sits exactly on the constraint, as the solution of a chance-constrained
program must. (b) Cleanup time: $0$\,d for 39/40 scenarios, the horizon cap for the
single violating one. Generated by \code{SCRIPTS/make\_figures.py}.}
\label{fig:chance}
\end{figure}

\subsection{Three dimensions: verified first, and that is how the bug was found}
\label{sec:res3d}

The 3D runner is a self-contained NumPy/SciPy implementation (TPFA Darcy,
diagonal Bear dispersion, vertical anisotropy, partially penetrating wells) with
pattern search and multiple starts. Before quoting any of its numbers we asked
it questions whose answers do not depend on our own code, and that turned out to
matter more than expected.

Two checks are now permanent tests. The first is a \emph{moment} test: put a
pulse in the middle of the domain and read the growth of its spatial variance.
It must obey $\mathrm{Var} = 2Dt$ with the coefficients the operator claims to
assemble, and it does --- longitudinal $D = \alpha_L v + v\Delta x/2$ (the
second term is the first-order upwind contribution, recovered to $10\%$),
transverse $D = \alpha_T v$ on both axes, and a centroid drifting at the pore
velocity $v = K\,\mathrm{grad}/\phi$ rather than the Darcy velocity. The second
is a \emph{front} test against Ogata--Banks. A mass-flux inlet is not a Dirichlet
inlet: the classical formula carries a reflected term
$\mathrm{e}^{vx/D}\operatorname{erfc} b$ that a purely advective inlet has no
analogue for, and at the parameters used here the two solutions differ by a
\emph{shift} of $D/v$ --- five to ten cells, independent of $\Delta x$. Against
the correct reference, the swept front
$\tfrac12 C_0\operatorname{erfc}\big((x-vt)/2\sqrt{Dt}\big)$, the interior
relative $L^2$ error is $1.1\times10^{-2}$ at $\Delta x = 1$\,m and
$4.0\times10^{-3}$ at $\Delta x = 0.5$\,m: first order in $\Delta x$, inside the
$1.5\times10^{-2}$ gate. The $0.41$ ``shape error'' we had been carrying was the
wrong reference, and no amplitude fit could have revealed that, because the
mismatch is a shift rather than a scale.

Writing those tests is what exposed a real defect. The 3D fields are stored
$(n_z, n_y, n_x)$, but \code{Grid3D.centres()} returned
\code{meshgrid(x, y, z)} --- that is, $(n_x, n_y, n_z)$ --- and the Gaussian well
kernel is built from those arrays and then flattened. Every well was therefore
written into a transposed cell: on a cubic check grid a well requested at
$(3.5, 2.5, 2.0)$ landed at $(0.5, 2.5, 3.5)$. No run could reveal it, because
the coordinates reported in \code{metrics.json} are the \emph{requested} ones and
the z-screen mask was applied to a coordinate array that was itself transposed.
Every 3D result before 2026-11-23 describes a well layout nobody specified. The
fix is three lines and two regression tests that ask where the kernel actually
peaks; the corrected results are in Section~\ref{sec:res3dcert}. We report the
episode because it is the argument for the protocol rather than against it: the
audit could not have found this, and neither could a parity test between the
NumPy and torch paths, which agreed with each other because both were wrong in
the same way. Only an outside reference could.

\begin{table}[htbp]
\centering
\small
\caption{3D presets after the well-placement fix. The paper preset now settles
on two wells, not four, and is $18\%$ cheaper.}
\label{tab:3d}
\begin{tabular}{lrrrr}
\toprule
preset & grid & $n_{KL}$ & scen.\ (tr/te) & result \\
\midrule
smoke & $12\cdot6\cdot3=216$ & 164 & 6/12 & degenerate: already compliant, $q^\star=0$ \\
paper & $32\cdot16\cdot8=4096$ & 1419 & 64/128 & $q^\star=[0,10.64,0,0,0,12.79,0,0,0]$, rel.\ $0.953$ \\
paper & $32\cdot16\cdot8=4096$ & 1419 & 64/128 & train $J=449{,}972$; certified $0.9725$ on 400 fresh fields \\
\bottomrule
\end{tabular}
\end{table}

\begin{figure}[htbp]
\centering
\includegraphics[width=0.46\textwidth]{figs/3d_well_layout.png}\hfill
\includegraphics[width=0.50\textwidth]{figs/3d_convergence.png}
\caption{3D paper preset after the fix: well layout and search convergence.
Artifacts: \code{out3d\_colab/paper\_fixed/figures/}.}
\label{fig:3d}
\end{figure}

\subsection{Two certified 3D policies, and what multiplicity looks like}
\label{sec:res3dcert}

The paper preset resolves the $20$\,m correlation length with $2.4$ cells, so a
second 3D geometry was run at $40\times20\times10=8{,}000$ cells
($\Delta x=\Delta y=5$\,m, $\Delta z=2$\,m, $n_{KL}=2062$), where the same
correlation length spans four cells and the screened interval covers three
layers, so partially penetrating wells produce real vertical flow. Every policy
below was audited on 400 fresh fields (seed base $900{,}000$) that no search had
seen.

\begin{table}[htbp]
\centering
\footnotesize
\caption{Corrected 3D results. Every row is a different search, and each was
audited on 400 fresh fields. The two $0.97$ rows are certified; the $M_{\text{train}}
=24$ row is the cheapest and inconclusive, and the NumPy row has the best in-run
number of any 3D run and still does not certify.}
\label{tab:3dcert}
\begin{tabular}{llrrrr}
\toprule
geometry & search & $q^\star$ (m$^3$/d) & train $J$ & in-run & 400-field rel.\;[CI] \\
\midrule
$8{,}000$ & torch T4, $M{=}24$  & $[0,0,12.5,0,0,8.59]$ & $441{,}298$ & $0.958$ & $0.9375$;\;[0.909, 0.957] \\
$8{,}000$ & torch T4, $M{=}40$  & $[0.05,12.5,0,0,10.06,0]$ & $446{,}742$ & $0.979$ & $0.9700$;\;[0.948, 0.983] \\
$8{,}000$ & torch T4, $M{=}100$ & $[0.2,12.5,0,0,9.77,0]$ & $446{,}286$ & $0.979$ & $0.9700$;\;[0.948, 0.983] \\
$4{,}096$ & torch T4, $M{=}64$  & $[0,10.64,0,0,0,12.79,0,0,0]$ & $449{,}972$ & $0.953$ & $0.9725$;\;[0.951, 0.985] \\
$4{,}096$ & NumPy, 3 starts    & $[17.05,0,0,0,0,0,0,0,8.42]$ & $457{,}397$ & $0.984$ & $0.9475$;\;[0.921, 0.965] \\
$13{,}824$ & torch T4, $M{=}40$ & $[12.7,0,0,0,0,9.28]$ & $445{,}090$ & $0.938$ & $0.9500$;\;[0.924, 0.967] \\
\bottomrule
\end{tabular}
\end{table}

Three things follow, and the first corrects an earlier draft of this section.
\emph{The $M_{\text{train}}=40$ failure was an artefact of the transposed wells.}
Before the fix, that search found a cheaper policy which audited at $0.9050$ and
looked like the two-dimensional story repeating. At the corrected geometry the
same search finds the $M_{\text{train}}=100$ policy to $0.1\%$ and both certify
at $0.970$; the reading of the earlier failure is withdrawn. A weaker boundary
does remain: $M_{\text{train}}=24$ ($\alpha M_{\text{train}}=1.2$) is the cheapest
policy in the table, $21.09$\,m$^3$/d at $J=441{,}298$, and it audits at
$0.9375$ with a point estimate under the gate, while $\alpha M_{\text{train}}=2.0$
passes. So cost and reliability are anti-correlated across these policies ---
$1.1\%$ cheaper, $0.037$ lower on the audit --- but the crossover sits between
one and two expected tail samples, lower than the two-dimensional rule of five.
We report the boundary we measured rather than the one we expected.

Second, \emph{policy multiplicity is real and is not a search artifact}. At the
paper geometry two searches with the same seeds and different hardware settled on
different designs: wells $\{2,6\}$ at $10.64$ and $12.79$\,m$^3$/d ($J=449{,}972$,
audited $0.9725$), and wells $\{1,9\}$ at $17.05$ and $8.42$\,m$^3$/d
($J=457{,}397$). Both are plausible engineering answers, $1.6\%$ apart in cost. A
single search is therefore not evidence of a unique optimum, and we do not claim
one.

Third, \emph{the in-run test row is not a certificate}, and the paper geometry
makes that unambiguous. The second design above has the \emph{best} in-run
reliability of any 3D run here --- $0.984$, comfortably above the gate --- and its
$400$-field audit lands at $0.9475$, inconclusive, with a point estimate under the
gate. Had we stopped at the run's own row we would have certified the wrong
policy; the audit reversed the ranking relative to the in-run numbers
($0.984$ versus $0.953$ in-run, $0.9475$ versus $0.9725$ audited).

A larger geometry does not extend either conclusion. At $13{,}824$ cells ($\Delta x=\Delta y=4.17$\,m, $n_{KL}=2706$,
$M_{\text{train}}=40$, $5.8$\,s per evaluation on the same T4, $103$ evaluations in
$613$\,s), extends neither conclusion. It found
$q^\star=[12.7,0,0,0,0,9.28]$\,m$^3$/d at $J=445{,}090$, and the audit returns
$J=445{,}101$ --- the objective reproduces across draws to $0.02\%$ --- with
$20/400$ violations, reliability $0.9500$ and Wilson $[0.924,0.967]$:
inconclusive, point estimate above the gate and lower bound below it. Its in-run
row had already failed ($0.9375$), so on this geometry the smoke detector was
right, which is the reverse of the paper geometry and is exactly why it is only a
smoke detector. What the run does \emph{not} do is transfer the $8{,}000$-cell
certificate ($0.9700$): the grid and the KL realisation change together
($n_{KL}$ 2062 versus 2706), so this is not a resolution study and we do not read
it as one. Cost per scenario keeps growing with cell count --- $0.032$\,s at
$4{,}096$, $0.071$ at $8{,}000$, $0.145$ at $13{,}824$ --- and the run still fits
a twelve-minute VM window.

What made the $100$-scenario search affordable is exactly the lever of
Section~\ref{sec:rescost}: the batched backend evaluates $100$ scenarios in
$7.2$\,s where the NumPy reference needs about $105$\,s, a $15\times$ gain from
replacing $100$ sparse factorisations with one batched Krylov solve. It is not a
property of the card --- CPU torch is already $2.4\times$ faster than NumPy at
this geometry --- and the T4 run cost $697$\,s in total.

\subsection{Cost: the accelerator is not the lever}
\label{sec:rescost}

\begin{table}[htbp]
\centering
\small
\caption{Measured cost levers. The first two are real and reproducible; the
third is a caution.}
\label{tab:cost}
\begin{tabularx}{\textwidth}{@{}lXr@{}}
\toprule
Lever & Measurement & Speedup \\
\midrule
Batch a compass sweep & one kernel call per sweep at $80\times40$, T4, $M=8$ & $47\times$ \\
Mixed float32 & float32 operator, float64 Krylov recurrence, $160\times96$ & $2.46\times$ \\
A100 vs T4 & same geometry; fixed launch overhead dominates & no gain \\
\bottomrule
\end{tabularx}
\end{table}

Table~\ref{tab:cost} is the cost section in one line: batching and precision
are levers, the accelerator name is not. Section~\ref{sec:res3dcert} is what
that advice buys when it is followed: batching a 3D sweep and its scenarios is
what made a $100$-scenario search affordable, and the affordable search is the
one that produced a policy that passes its audit. The workload is latency-bound in
float32 and fp64-throughput-bound in float64, so the choice of accelerator is
regime-dependent rather than monotone in price. The honest conclusion is the
one in \code{DOCS/GPU.md}: do not rent a GPU for this problem until the kernel
launch count is reduced; the measured count is $4{,}022$ matvecs per evaluation,
$1.7\times$ the number the optimization plan assumed. Figure~\ref{fig:compute}
shows the same table as bars, grouped by geometry so that no bar is compared
against a different grid.

\begin{figure}[htbp]
\centering
\includegraphics[width=0.92\textwidth]{figs/compute_cost.png}
\caption{Measured cost, grouped so that every comparison is like-for-like.
(a) One 18-candidate compass sweep. At $80\times40$ float32 the T4 is $47\times$
faster than one NumPy core; at $160\times96$ float64 it is \emph{not faster at
all} ($680$\,s against $662$\,s), and only the A100 wins ($106$\,s, $6.4\times$).
Float32 is launch-overhead-bound and float64 is fp64-throughput-bound, so
``buy an accelerator'' is not one decision. (b) On the T4 the $17\times$ comes
from batching candidates into one kernel call, not from the card. Numbers
measured on rented hardware; see \code{DOCS/GPU.md}.}
\label{fig:compute}
\end{figure}
